\documentclass[12pt, oneside]{article}

\usepackage{xparse}

\usepackage{geometry}
\geometry{letterpaper}
\usepackage[parfill]{parskip}
\usepackage{float}

% Math packages
\usepackage{amssymb}
\usepackage{amsmath}
\usepackage{amsthm}
\usepackage{mathalpha}

% Tables
\usepackage{booktabs}
\usepackage{afterpage}

% Hyperlinks
\usepackage{hyperref}

% References
\numberwithin{equation}{section}
\usepackage[numbers]{natbib}

% TikZ (optional for later)
\usepackage{tikz}
\usetikzlibrary{shapes,arrows,chains}

% Theorem environments
\theoremstyle{definition}
\newtheorem{definition}{Definition}[section]
\newtheorem{proposition}{Proposition}[section]
\newtheorem{axiom}{Axiom}[section]

\theoremstyle{plain}
\newtheorem{theorem}{Theorem}[section]
\newtheorem{lemma}[theorem]{Lemma}
\newtheorem{corollary}[theorem]{Corollary}
\newtheorem{example}[theorem]{Example}

\theoremstyle{remark}
\newtheorem*{remark}{Remark}

\title{Affine Constraint Dynamics and Density-One Descent in the Collatz Map}
\author{Wayne Brassem}

\begin{document}

% Custom commands
\NewDocumentCommand{\setN}{}{\mathbb{N}}
\NewDocumentCommand{\setZ}{}{\mathbb{Z}}
\NewDocumentCommand{\setQ}{}{\mathbb{Q}}

% Document commands which provide hyperlinks to OEIS sequences referenced herein
\NewDocumentCommand{\OEIS}{m}{\href{https://oeis.org/#1}{#1}}

\maketitle

\begin{abstract}

We study the Collatz map via the fully accelerated transformation
\[
f(x)=\frac{3x+1}{2^{v_2(3x+1)}}.
\]

Finite trajectories are encoded by division words, each inducing an affine map
whose admissibility imposes congruence constraints on initial values. This
yields a symbolic–affine framework in which admissible words correspond to
residue classes modulo powers of two defined by these constraints.

We analyze the growth of admissible words through their associated division
counts, showing that the induced affine constraints refine at an exponential
rate governed by $2^{D(\omega)} \sim 3^m$, while the number of admissible words
grows subcritically like $\rho^m$ for some $\rho < 3$. This imbalance forces
an exponential sparsification of admissible initial conditions.

As a consequence, the set of integers that fail to admit a finite descent prefix
has natural density zero, and therefore almost all integers eventually decrease
under iteration of the Collatz map.

\end{abstract}

\newpage
\tableofcontents

% ============================================================
\newpage
\section{Introduction}

The $3x+1$ problem, also known as the Collatz conjecture~\cite{Collatz1937}, concerns
the behavior of the map
\[
T(x) =
\begin{cases}
\displaystyle
    \phantom{3x} \frac{x}{2} & \text {if } x \text{ is even}, \\[8pt]
\displaystyle
    3x+1,                    & \text{if $x$ is odd}.
\end{cases}
\]
on the positive integers. It asserts that every positive integer eventually reaches $1$
under iteration of $T$. Despite its elementary formulation, the conjecture has resisted
proof for decades~\cite{Lagarias2010,Terras1976}. Extensive computational verification
confirms convergence for very large ranges~\cite{OliveiraESilva2010}, but does not
reveal the global structure governing trajectory behavior.

A major advance by Tao~\cite{Tao2019} showed that almost all integers, in the sense of
logarithmic density, eventually decrease under iteration of the Collatz map. This result
is obtained via probabilistic averaging rather than an explicit structural decomposition
of the integers.

The present work develops a deterministic, arithmetic framework for analyzing typical
Collatz behavior. Rather than studying individual trajectories, we encode the dynamics
symbolically using the fully accelerated map
\[
f(x)=\frac{3x+1}{2^{v_2(3x+1)}},
\]
which maps odd integers to odd integers by compressing each odd step into a single
transformation.

Each trajectory under $f$ is represented by a finite \emph{division word}, recording the
2-adic valuation at each iterate. These words admit an exact affine representation:
each division word induces a map of the form
\[
F_\omega(x) = \frac{3^m x + c(\omega)}{2^{D(\omega)}},
\]
and determines a corresponding congruence class modulo $2^{D(\omega)}$.
Thus admissible division words determine canonical arithmetic
\emph{constraint fibers} in the integers.

This correspondence allows the dynamics to be studied at the level of residue classes
rather than individual orbits. Convergent division words—those for which the associated
affine map is contractive—identify families of integers that decrease after a fixed
number of steps. The global problem is thereby reduced to understanding how these
families are distributed and how their densities accumulate.

Two competing effects govern this distribution. On one hand, the number of admissible
convergent division words grows exponentially with length. On the other hand, each word
imposes a congruence constraint whose modulus grows exponentially with the total 2-adic
valuation. Using structural constraints on admissible division words, encoded by an
increment grammar, we show that the combinatorial growth rate is strictly subcritical:
the number of admissible words grows like $\rho^m$ for some $\rho < 3$, while the
moduli scale like $2^{D(\omega)} \sim 3^m$.

This transforms the Collatz problem into a competition between symbolic growth
and arithmetic refinement, with convergence governed by which process dominates.

This imbalance leads to \emph{exponential sparsification}: the proportion of integers
admitting a convergent division word of length $m$ decays exponentially in $m$.
Summing over all lengths shows that the complement—integers that avoid all such
descent patterns—has natural density zero.

Consequently, the set of integers whose trajectory under iteration of $f$ eventually
decreases has natural density one.

While this does not exclude the existence of exceptional trajectories, any such
exceptions must lie within a set of zero density and cannot arise from any scalable
arithmetic structure generated by the symbolic dynamics.

\paragraph{Structure of the Paper.}
The argument proceeds as follows:
\begin{itemize}
\item Section~\ref{sec:division-counts} introduces division words and their total division counts.
\item Section~\ref{sec:affine-representation} develops the affine representation of division words and their congruence structure.
\item Section~\ref{sec:affine-grammar} establishes admissibility and minimal valuation constraints.
\item Section~\ref{sec:numerator-growth} controls the growth of admissible words via the increment grammar.
\item Section~\ref{sec:affine-constraint-fibers} identifies the residue class structure induced by division words.
\item Section~\ref{sec:exponential-sparsification} proves exponential sparsification of admissible sets.
\item Section~\ref{sec:density-one-descent} establishes density-one descent.
\end{itemize}

% ============================================================
\section{Symbolic Encoding of Collatz Trajectories}
\label{sec:division-counts}

% ------------------------------------------------------------
\subsection{Division Words}

\begin{definition}[Division Word]
A division word of length $m \ge 0$ is a sequence
\[
\omega = (d_0, \dots, d_{m-1}),
\]
where each $d_i \in \mathbb{N}$ represents the number of factors of $2$
removed in the $i$-th iterate of the accelerated map
\[
f(x) = \frac{3x+1}{2^{v_2(3x+1)}}.
\]
\end{definition}

\begin{remark}
Since $f$ maps odd integers to odd integers, each $3f^i(x)+1$ is even,
so $d_i \ge 1$ for all $i$.
\end{remark}

\begin{definition}[Realization]
An integer $x \in \mathbb{N}$ realizes $\omega$ if
\[
d_i = v_2(3f^i(x)+1), \quad 0 \le i < m.
\]
\end{definition}

% ------------------------------------------------------------
\subsection{Total Division Count}

\begin{definition}
The total division count of a word $\omega$ is
\[
D(\omega) := \sum_{i=0}^{m-1} d_i.
\]
\end{definition}

\begin{remark}
The quantity $D(\omega)$ measures the total dyadic contraction
associated with the symbolic path $\omega$.
\end{remark}

% ============================================================
\section{Affine Representation of Division Words}
\label{sec:affine-representation}

\begin{definition}
Each division word $\omega$ of length $m$ induces a formal affine map
\[
F_\omega(x) = \frac{3^m x + c(\omega)}{2^{D(\omega)}},
\]
obtained by symbolic composition of the accelerated map.
\end{definition}

\begin{proposition}
The map $\omega \mapsto F_\omega$ defines a semigroup homomorphism
from concatenation of division words to composition of affine maps.
\end{proposition}

% ------------------------------------------------------------
\subsection{Example}

Consider convergent segments with $m=4$ iterates of $f$.  
The minimal total division count for $m=4$ is
\[
A020914(4) = 7.
\]

One admissible division word attaining this minimum is
\[
\omega = (1,1,1,4),
\]
for which
\[
D(\omega) = 1+1+1+4 = 7.
\]

Composing the accelerated map symbolically yields
\[
f_\omega(x)=\frac{3^4 x + c(\omega)}{2^7}
           = \frac{81x + 65}{128},
\]
so the canonical affine constant for this word is
\[
c(\omega)=65.
\]

The integrality condition
\[
81x + 65 \equiv 0 \pmod{128}
\]
determines a unique residue class
\[
x \equiv 15 \pmod{128}.
\]
This congruence defines a constraint fiber with modulus $2^7$.

\medskip

\noindent
\textbf{Interpretation.}
The word $\omega = (1,1,1,4)$ means:

\begin{itemize}
    \item At each of the first three iterates of $f$, $3x+1$ contains exactly one factor of $2$.
    \item At the fourth iterate of $f$, $3x+1$ contains four factors of $2$.
\end{itemize}

Under the accelerated map
\[
f(x)=\frac{3x+1}{2^{v_2(3x+1)}},
\]
all of these divisions occur within a single iterate. Thus $d_3=4$ means that
the fourth iterate divides by $2^4$ in that step. The constant $c(\omega)$ is
fixed entirely by this symbolic structure and ensures that precisely the
integers in the residue class $x \equiv 15 \pmod{128}$ produce integer iterates
under the affine form.

\medskip

This example illustrates that the behavior of trajectories is governed
entirely by the symbolic data encoded in $\omega$ and its associated
affine map. The admissible initial values are precisely those satisfying
the induced congruence constraint, independent of any particular choice
of representative.

% ============================================================
\section{Affine Symbolic Admissibility Grammar}
\label{sec:affine-grammar}

\begin{remark}[Division Grammar as a Dynamical Encoding]
Recall from Section~\ref{sec:division-counts} that a \emph{division word}
$\omega=(d_0,\dots,d_{m-1})$ records the exact 2-adic valuation
$d_i = v_2(3 x_i + 1)$, occurring at each iterate of the fully accelerated
Collatz map.

Unlike a purely combinatorial encoding, division words retain essential
dynamical information. Each division word $\omega=(d_0,\dots,d_{m-1})$
determines the sequence of multiplicative factors of the linear coefficient
\[
r_i = \frac{3}{2^{d_i}},
\]
so that the contractive or expansive contribution of each iterate is
preserved in the linear coefficient (with the affine term handled separately
via $c(\omega)$). This structure enables quantitative control of admissible
trajectories via affine growth and to formulate sufficient conditions for
descent, even though individual trajectories are collapsed under the encoding.
\end{remark}

% ------------------------------------------------------------
\subsection{Explicit Form of the Affine Constant}
\begin{lemma}[Explicit Form of the Affine Constant]
Let $\omega = (d_0,\dots,d_{m-1})$ be an admissible division word.  
Then the affine map
\[
F_\omega(x) = \frac{3^m x + c(\omega)}{2^{D(\omega)}}
\]
has the explicit form
\[
c(\omega)
=
\sum_{j=0}^{m-1}
3^{m-1-j}\cdot 2^{\sum_{i=0}^{j-1} d_i},
\]
where empty sums are defined as $0$.
\end{lemma}

% ------------------------------------------------------------
\subsection{Injectivity of the Affine Encoding}
\begin{lemma}[Injectivity of the Affine Encoding at Fixed Length]
\label{lem:injectivity}
Let $\omega = (d_0,\dots,d_{m-1})$ and $\omega' = (d'_0,\dots,d'_{m-1})$
be admissible division words of the same length $m$. Let
\[
c(\omega)
=
\sum_{j=0}^{m-1}
3^{m-1-j}\cdot 2^{D_j},
\quad
D_j = \sum_{i=0}^{j-1} d_i,
\]
and define $c(\omega')$ analogously.

If $\omega \ne \omega'$, then
\[
c(\omega) \ne c(\omega').
\]

Equivalently, the map $\omega \mapsto c(\omega)$ is injective on admissible
division words of fixed length $m$.
\end{lemma}

\begin{proof}
Suppose $\omega \ne \omega'$, and let
\[
k = \min\{\, j : d_j \ne d'_j \,\}
\]
be the first index at which the two words differ. Then for all $j < k$,
we have $D_j = D'_j$, while $D_k \ne D'_k$.

Consider the difference
\[
\Delta := c(\omega) - c(\omega')
=
\sum_{j=0}^{m-1}
3^{m-1-j}\bigl(2^{D_j} - 2^{D'_j}\bigr).
\]
The terms with $j<k$ cancel, so
\[
\Delta
=
3^{m-1-k}\bigl(2^{D_k} - 2^{D'_k}\bigr)
+
\sum_{j=k+1}^{m-1}
3^{m-1-j}\bigl(2^{D_j} - 2^{D'_j}\bigr).
\]

Let $\delta_k = \min(D_k, D'_k)$. Then
\[
2^{D_k} - 2^{D'_k}
=
2^{\delta_k}(2^a - 2^b),
\]
where one of $a,b$ is zero and $2^a - 2^b$ is odd. Hence
\[
v_2\!\left(3^{m-1-k}\bigl(2^{D_k} - 2^{D'_k}\bigr)\right)
=
\delta_k,
\]
since $3^{m-1-k}$ is odd.

For $j>k$, both $D_j$ and $D'_j$ strictly exceed $\delta_k$, since the
prefix sums increase by positive integers. Therefore
\[
v_2\!\left(3^{m-1-j}\bigl(2^{D_j} - 2^{D'_j}\bigr)\right)
>
\delta_k
\quad \text{for all } j>k.
\]

It follows that $\Delta$ is a sum of one term with 2-adic valuation
exactly $\delta_k$ and remaining terms of strictly higher 2-adic
valuation. Hence
\[
v_2(\Delta) = \delta_k,
\]
and in particular $\Delta \ne 0$.

Therefore $c(\omega) \ne c(\omega')$, completing the proof.
\end{proof}

This implies distinct division words induce distinct affine constraints.

% ------------------------------------------------------------
\subsection{Stepwise Realizability}

\begin{theorem}[Stepwise Realizability of Division Words]
\label{thm:stepwise-realizability}
Let $\omega = (d_0,\dots,d_{m-1})$ be a division word, and define $c(\omega)$ by
\[
c(\omega)
=
\sum_{j=0}^{m-1}
3^{m-1-j}\cdot 2^{\sum_{i=0}^{j-1} d_i}.
\]

Then the identity
\[
2^{D(\omega)} x_m = 3^m x_0 + c(\omega)
\]
holds if and only if the sequence $(x_0,\dots,x_m)$ satisfies the accelerated
Collatz relations
\[
x_{i+1} = \frac{3x_i+1}{2^{d_i}}, \quad 0 \le i < m,
\]
with each intermediate step integral.

In particular, the congruence
\[
3^m x_0 + c(\omega) \equiv 0 \pmod{2^{D(\omega)}}
\]
is equivalent to the existence of a trajectory realizing $\omega$.
\end{theorem}

\begin{remark}[Structural Interpretation]
Each term in $c(\omega)$ corresponds to the additive contribution of a single
``+1'' in the Collatz iteration, propagated forward by subsequent multiplications
by $3$ and backward-scaled by preceding divisions by $2$.
\end{remark}

% ------------------------------------------------------------
\subsection{Admissibility of Division Words}

\begin{definition}[Admissibility]
\label{def:admissible-division-word}
A division word $\omega$ is admissible if the congruence
\[
3^m x + c(\omega) \equiv 0 \pmod{2^{D(\omega)}}
\]
has a solution.
\end{definition}
\begin{remark}
Admissibility is therefore an arithmetic condition on the word $\omega$
itself, expressed as a divisibility constraint on the associated affine
map. When this condition holds, the set of all integers $x$ for which
$F_\omega(x)$ is integral defines a unique congruence constraint modulo
$2^{D(\omega)}$, corresponding to the residue class induced by $\omega$.
\end{remark}

% ------------------------------------------------------------
\subsection{Minimal Division Count}

\begin{lemma}[Minimal Division Count]
For any admissible division word $\omega=(d_0,\dots,d_{m-1})$ of length $m$,
\[
d_0 + \cdots + d_{m-1} \;\ge\; A020914(m),
\]
where
\[
A020914(m) = \lceil m \log_2(3) \rceil.
\]
Equality holds if and only if $\omega$ achieves the minimal possible total
division count among admissible words of length $m$.
\end{lemma}

% ------------------------------------------------------------
\subsection{Convergent Division Word}

\begin{definition}[Convergent Division Word]
An admissible division word $\omega = (d_0,\dots,d_{m-1})$ is said to be
\emph{convergent} (in the linear sense) if
\[
\frac{3^m}{2^{d_0+\cdots+d_{m-1}}} < 1.
\]
\end{definition}

% ============================================================
\section{Numerator Growth and Counting Control}
\label{sec:numerator-growth}

\begin{proposition}
The number of convergent division words of length $m$ is given by \OEIS{A186009}.
This sequence therefore describes the combinatorial growth rate of convergent
division words at depth $m$, and hence of the associated residue classes.
\end{proposition}

\begin{remark}[Canonical Tabular Realization]
A canonical realization of this counting sequence via the horizontal-sum construction,
together with its derivation from the increment grammar induced by \OEIS{A022921}, is given
in Appendix~\ref{app:spine-growth}. The appendix serves as a structural reference:
the arguments in this section rely only on the grammar itself and not on any numerical
properties of the tabulation.
\end{remark}

\begin{remark}[Grammar Suppression Mechanism]
\label{rem:grammar-suppression}
The increment grammar prohibits arbitrarily long sequences of locally expansive steps.
Any such expansion must be separated by configurations that reduce the horizontal sum.
\end{remark}

% ------------------------------------------------------------
\subsection{Controlled Growth of Horizontal Sums}

\begin{lemma}[Controlled Growth of Horizontal Sums]
\label{lem:controlled-horizontal-growth}

Let a row have strictly positive entries with total sum $S$.

A single increment step produces a new row whose total sum is $S$.

If a double increment is applied immediately afterward, producing a row in which
the terminal entry equals the sum of the previous row and all other entries remain
strictly positive, then the total sum of the resulting row satisfies
\[
S_{\mathrm{new}} < 3S.
\]
\end{lemma}

\begin{proof}
The total sum of the original row is $S$.

Under a double increment step, the resulting row consists of:
\begin{enumerate}
    \item (i) interior entries with total $S_{\mathrm{int}} < S$, and
    \item (ii) two terminal contributions, each equal to $S$.
\end{enumerate}

Hence the total sum satisfies
\[
S_{\mathrm{new}} = S_{\mathrm{int}} + 2S < 3S.
\]
\end{proof}

% ------------------------------------------------------------
\subsection{Dyadic–Ternary Comparison}
\begin{lemma}[Dyadic–Ternary Comparison]
\label{lem:uniform-control}
For all $i \ge 0$,
\[
3^i < 2^{A020914(i)} < 2\cdot 3^i.
\]
Consequently, the density contribution of convergent residue classes of
stopping time $i$, given by \(\widetilde A(i) 2^{-A020914(i)}\), satisfies
\[
\frac{\widetilde A(i)}{2 \cdot 3^i}
\;<\;
\widetilde A(i)\,2^{-A020914(i)}
\;<\;
\frac{\widetilde A(i)}{3^i}.
\]
\end{lemma}

\begin{proof}
By definition, $A020914(i)$ is the minimal integer $m$ such that $2^m > 3^i$.
Consecutive powers of two differ by a factor of $2$, so
\[
3^i < 2^{A020914(i)} < 2 \cdot 3^i.
\]
Taking reciprocals and multiplying by $\widetilde A(i)$ gives the stated bounds.
\end{proof}

% ------------------------------------------------------------
\subsection{Maximal Admissible Growth Rate}
\begin{theorem}[Maximal Admissible Growth Rate]
\label{thm:maximal-growth}
Let $\widetilde A(i) = A186009(i+1)$ denote the number of convergent residue classes
of stopping time $i$, generated by the horizontal-sum construction associated with
the increment grammar of \OEIS{A022921}.  Let $B(i) := A020914(i)$ denote the
cumulative exponent of $2$ up to level $i$.  

Then there exists a constant $\rho < 3$ such that
\[
\widetilde A(i) \le C \rho^i,
\]
and consequently
\[
\frac{\widetilde A(i)}{2^{B(i)}} \to 0.
\]
\end{theorem}

\begin{proof}
The increment sequence \OEIS{A022921} induces a finite symbolic grammar on $\{s,d\}$,
corresponding to single and double increments in the horizontal-sum construction.

A detailed analysis of the increment grammar (see Appendix~\ref{app:spine-grammar})
shows that all admissible increment sequences have asymptotic growth rate bounded by a
uniform constant $\rho < 3$.

We denote the optimal such bound by $\rho_5 < 3$.

Any admissible infinite increment sequence decomposes into concatenations of finite
blocks permitted by the grammar. As explained in Remark~\ref{rem:grammar-suppression},
all locally expansive blocks are separated by mandatory contractive buffering. Horizontal-sum
growth is submultiplicative under concatenation, so the asymptotic per-step geometric
growth is bounded above by the maximal geometric mean among admissible blocks, which is
attained by the $5$-cycle and denoted $\rho_5$. Hence
\[
\limsup_{i\to\infty} \widetilde A(i)^{1/i} \le \rho_5 < 3.
\]

From Lemma~\ref{lem:uniform-control}, the denominator satisfies
\[
3^i < 2^{B(i)} < 2\cdot 3^i.
\]
Consequently, there exists a constant $C>0$, depending only on initial conditions
and finite prefixes, such that
\[
\widetilde A(i) \le C\,\rho_5^i \quad \text{for all } i.
\]
Dividing by $2^{B(i)}$ gives
\[
\frac{\widetilde A(i)}{2^{B(i)}} \le \frac{C\,\rho_5^i}{3^i} \longrightarrow 0,
\]
which completes the proof.
\end{proof}

Thus the exponential growth of admissible division words is strictly dominated
by the exponential growth of the denominator $2^{B(i)} \sim 3^i$.

% ------------------------------------------------------------
\subsection{Spectral Bracketing of Symbolic Growth}
\label{sec:spectral-bracketing}
The symbolic growth of admissible configurations is governed by two competing
envelope systems.

On the lower end, the restricted grammar underlying \OEIS{A047749} forbids
consecutive doubling events and yields an explicitly solvable contraction-limited
model with exponential rate
\[
\lambda_{\mathrm{low}} = \frac{3\sqrt{3}}{2}.
\]

This is explained in detail in Appendix~\ref{app:A047749-lower-bound}, where the
growth rate is derived from the dominant eigenvalue of the associated transition matrix.

On the upper end, the full admissible grammar of \OEIS{A186009} is controlled by
cycle decomposition, with maximal spectral growth rate
\[
\lambda_{\mathrm{high}} = \rho_5 < 3.
\]

Thus the full symbolic system is spectrally bracketed:
\[
\lambda_{\mathrm{low}} \;\le\; \text{growth rate of admissible structure} \;\le\; \lambda_{\mathrm{high}}.
\]

% ============================================================
\section{Affine Constraint Fibers}
\label{sec:affine-constraint-fibers}

% ------------------------------------------------------------
\subsection{Congruence Stability of Division Words}

\begin{theorem}[Congruence Stability of Division Words]
Let $\omega$ be admissible and let $r_\omega$ satisfy
\[
r_\omega \equiv -c(\omega)\cdot (3^m)^{-1} \pmod{2^{D(\omega)}}.
\]
Then for all $x \equiv r_\omega \pmod{2^{D(\omega)}}$, the first $m$
iterations of the accelerated Collatz map realize exactly the division word $\omega$.
\end{theorem}

% ------------------------------------------------------------
\subsection{Residue Class Representation of Division Words}

\begin{theorem}[Affine Constraint Representation]
Let $\omega$ be an admissible division word of length $m$.

Then there exists a unique residue $r_\omega \pmod{2^{D(\omega)}}$ such that
\[
x \text{ realizes } \omega \iff x \equiv r_\omega \pmod{2^{D(\omega)}}.
\]

Thus admissible division words correspond to affine congruence constraints
of modulus $2^{D(\omega)}$.
\end{theorem}

\begin{proof}
From the affine representation,
\[
F_\omega(x) = \frac{3^m x + c(\omega)}{2^{D(\omega)}},
\]
integrality is equivalent to
\[
3^m x + c(\omega) \equiv 0 \pmod{2^{D(\omega)}}.
\]

Since $3^m$ is odd, it is invertible modulo $2^{D(\omega)}$, yielding
\[
x \equiv -c(\omega) \cdot (3^m)^{-1} \pmod{2^{D(\omega)}}.
\]

This determines a unique residue modulo $2^{D(\omega)}$.

Conversely, any $x$ in this class satisfies the integrality condition,
and therefore realizes the division word $\omega$.

Thus admissible division words are in one-to-one correspondence with affine
congruence constraints modulo powers of two.
\end{proof}

% ============================================================
\section{Exponential Sparsification}
\label{sec:exponential-sparsification}

\begin{lemma}[Constraint Refinement]
Each admissible division word $\omega$ imposes a congruence constraint
modulo $2^{D(\omega)}$, where $D(\omega) \ge A020914(m)$.
\end{lemma}

\begin{lemma}[Subcritical Symbolic Growth]
The number of admissible division words of length $m$ satisfies
\[
\#\{\omega : |\omega|=m\} = O(\rho^m)
\quad \text{for some } \rho < 3.
\]
\end{lemma}

\begin{proof}
This follows from Theorem~\ref{thm:maximal-growth}, which bounds the growth
of admissible configurations generated by the increment grammar.
\end{proof}

\begin{theorem}[Exponential Sparsification]
The set of integers admitting a division word of length $m$ occupies at most
a proportion on the order of
\[
\frac{\#\{\text{admissible words of length } m\}}{2^{A020914(m)}}.
\]
In particular, this proportion decays exponentially in $m$.
\end{theorem}

\begin{lemma}[Exponential Tail Decay]
Let $E_m$ denote the set of integers that do not admit a descent prefix
of length at most $m$.

Then $\mu(E_m)$ decays exponentially in $m$.
\end{lemma}

\begin{proof}
An integer fails to admit a descent prefix of length $m$ only if it avoids
all admissible convergent division words of length at most $m$.

Each such word imposes a congruence constraint modulo $2^{D(\omega)}$,
with $D(\omega) \ge A020914(m)$.

By subcritical symbolic growth, the number of such constraints grows
strictly slower than $2^{A020914(m)}$, so the total proportion of integers
covered by these constraints approaches $1$ exponentially fast.

Hence the complement $E_m$ decays exponentially.
\end{proof}

% ============================================================
\section{Density-One Descent}
\label{sec:density-one-descent}

\begin{lemma}[Exhaustion by Descent Prefixes]
Let $S_m$ denote the set of integers that admit no convergent division word
of length $\le m$. Then
\[
S_{m+1} \subseteq S_m,
\quad\text{and}\quad
\bigcap_{m=1}^\infty S_m
\]
is precisely the set of integers that admit no finite descent prefix.
\end{lemma}

\begin{theorem}[Density-One Descent]
The set of integers that do not admit any finite convergent division word
has natural density zero.
\end{theorem}

\begin{proof}[Proof of Theorem]
For each $m$, the complement of $S_m$ is the union of all residue classes
corresponding to convergent division words of length $\le m$.

Each convergent division word of length $m$ defines a residue class of
density $2^{-D(\omega)} \le 2^{-B(m)}$, and by subadditivity of natural density,
the total measure is boundedabove by the sum of their individual densities.

By the exponential sparsification result, the total number of such words
grows like $\widetilde A(m) \le C \rho^m$ with $\rho < 3$, while
$2^{B(m)} \sim 3^m$. Hence the total density contributed by words of length
$m$ decays exponentially.

Summing these exponentially decaying contributions over all lengths $\le m$
shows that the total density of integers admitting a convergent division word
of length at most $m$ tends to $1$ as $m \to \infty$.

Therefore
\[
\mu\!\left(\bigcap_{m=1}^\infty S_m\right) = 0,
\]
which proves the theorem.
\end{proof}

\begin{corollary}
The set of integers whose trajectory under $f$ eventually decreases
has natural density one.
\end{corollary}

\begin{remark}
Density emerges as a consequence of affine growth imbalance.
\end{remark}

% ------------------------------------------------------------
\subsection{Hypothetical Zero-Density Trajectories}

The density result established above does not exclude the logical possibility of
non-convergent Collatz trajectories. However, it severely restricts the form that
any such trajectory could take. In particular, any exceptional set must lie entirely
outside the union of convergent cylinders whose cumulative natural density
approaches one, confining any hypothetical exception to a set of natural density zero.

Any trajectory avoiding all convergent cylinders must eventually avoid every
admissible convergent division word, and hence cannot be captured by the cylinder
structures responsible for density accumulation. Such a trajectory, if it exists,
must therefore lie entirely within the residual set of density zero.

In this sense, any hypothetical exception would be \emph{dynamically isolated within
a set of zero natural density}.  It would not be supported by a scalable arithmetic
structure or by a cylinder hierarchy of positive density, but would instead persist
only as a sparse, non-propagating phenomenon within the integer space.

% ============================================================
\section{Discussion and Outlook}

% Interpretation of result
% Relation to classical approaches
% Remaining gap (zero-density exceptional set)
% Possible extensions

% ============================================================
% \newpage
% \appendix
% \section{Affine Constants and Cylinder Reconstruction}
% \label{app:affine-constants}

% Each division word $\omega = (d_0,\dots,d_{m-1})$ induces an affine map of length $m$
% (see Section~\ref{sec:division-counts})
% \[
% F_\omega(x) = \frac{3^m x + c(\omega)}{2^{D(\omega)}},
% \qquad
% D(\omega) = d_0 + \cdots + d_{m-1}.
% \]

% The constant $c(\omega)$ is uniquely determined modulo $2^{D(\omega)}$, and plays
% a central role in the induced congruence structure of the associated cylinder.

% This appendix collects three equivalent perspectives:
% (i) an explicit derivation of $c(\omega)$,
% (ii) the resulting residue class (cylinder) structure, and
% (iii) a computational reconstruction procedure.

% % ------------------------------------------------------------
% \subsection{Derivation of the Affine Constant}

% Let $\omega = (d_0,\dots,d_{m-1})$ and consider the accelerated iterations
% \[
% x_{i+1} = \frac{3x_i + 1}{2^{d_i}}, \quad 0 \le i < m.
% \]

% Repeated substitution yields the identity
% \[
% 2^{D(\omega)} x_m = 3^m x_0 + c(\omega),
% \]
% where $c(\omega)$ accumulates the contributions of the additive terms under
% forward multiplication by powers of $3$ and scaling by powers of $2$.

% Expanding inductively gives the explicit formula
% \[
% c(\omega)
% =
% \sum_{j=0}^{m-1}
% 3^{m-1-j}\cdot 2^{\sum_{i=0}^{j-1} d_i},
% \]
% with empty sums interpreted as $0$.

% This expression depends only on the division word $\omega$ and is independent
% of the initial representative.

% \begin{example}
% For $\omega = (1,1,1,4)$ with $m=4$ and $D(\omega)=7$, the expansion yields
% \[
% c(\omega) = 65,
% \]
% and hence
% \[
% F_\omega(x) = \frac{81x + 65}{128}.
% \]
% \end{example}

% % ------------------------------------------------------------
% \subsection{Cylinder (Residue Class) Structure}

% Admissibility of $\omega$ is equivalent to solvability of the congruence
% \[
% 3^m x + c(\omega) \equiv 0 \pmod{2^{D(\omega)}}.
% \]

% Since $3^m$ is invertible modulo $2^{D(\omega)}$, this determines a unique residue class
% \[
% x \equiv r_\omega \pmod{2^{D(\omega)}},
% \qquad
% r_\omega \equiv -c(\omega)\,(3^m)^{-1}.
% \]

% Thus all integers realizing $\omega$ are given by the cylinder
% \[
% x(k) = r_\omega + k\cdot 2^{D(\omega)}, \quad k \ge 0.
% \]

% Applying $F_\omega$ to this parametrization yields the affine form
% \[
% F_\omega(x(k)) = k\cdot 3^m + K_\omega,
% \]
% where
% \[
% K_\omega = \frac{3^m r_\omega + c(\omega)}{2^{D(\omega)}}.
% \]

% This identifies each division word with a full affine arithmetic progression in both
% domain and image.

% % ------------------------------------------------------------
% \subsection{Computational Recovery of $c(\omega)$}

% Assume that $\omega$ is admissible, so that the associated residue class modulo
% $2^{D(\omega)}$ is nonempty. Let $x$ be any integer realizing $\omega$.

% Then $c(\omega)$ can be recovered by:
% \begin{enumerate}
% \item Compute $y = 3^m x$.
% \item Let $r = \left\lceil \frac{y}{2^{D(\omega)}} \right\rceil$.
% \item Define
% \[
% c(\omega) = r\cdot 2^{D(\omega)} - y.
% \]
% \end{enumerate}

% This ensures
% \[
% 0 \le c(\omega) < 2^{D(\omega)}
% \]
% and reconstructs the unique representative of $c(\omega)$ modulo $2^{D(\omega)}$.

% \begin{remark}
% This computational procedure is not required for the structural or density arguments
% in the main text. It provides a concrete realization of the affine encoding from any
% representative of the corresponding cylinder.
% \end{remark}

% % ------------------------------------------------------------
% \subsection{Consistency}

% By the congruence characterization established in the main text, any two integers
% realizing the same division word differ by a multiple of $2^{D(\omega)}$.

% Hence the value of $c(\omega)$ is well-defined and depends only on $\omega$,
% not on the chosen representative.

% % % ============================================================
% % \newpage
% % \section{Extremal Growth Envelope}
% % \label{app:extremal-growth-envelope}

% % This section establishes the global growth bound governing admissible symbolic dynamics
% % under the increment grammar \OEIS{A022921}. It is the structural component used in
% % Theorem~\ref{thm:maximal-growth}.

% % % ------------------------------------------------------------
% % \subsection{Local Growth Structure}

% % The horizontal-sum construction decomposes admissible configurations into
% % a finite set of local patterns generated by single and double increment
% % steps (denoted $s$ and $d$). Each admissible pattern induces a
% % multiplicative change in total row sum, independent of global state.

% % \begin{lemma}[Local Growth Bounds]
% % \label{lem:local-growth-bounds}
% % Each admissible local pattern has bounded multiplicative growth:
% % \[
% % \begin{array}{c|c}
% % \text{Pattern} & \text{Max Growth Factor} \\
% % \hline
% % s,d & 3 \\
% % d,d & 123/33 \\
% % s,d,s & 341/131 \\
% % d,d,s & 329/123
% % \end{array}
% % \]
% % These bounds hold uniformly over all admissible placements in the horizontal-sum dynamics.
% % \end{lemma}

% % \begin{remark}
% % These constants arise from finite combinatorial enumeration of local
% % update rules. Their role in the global argument is solely to bound
% % per-step multiplicative growth; their precise derivation is not used
% % elsewhere in the main proof.
% % \end{remark}

% % % ------------------------------------------------------------
% % \subsection{Cycle Decomposition}

% % Admissible infinite words under the increment grammar decompose into
% % concatenations of finite cycles. Among these, a distinguished 5-cycle
% % acts as the extremal growth configuration.

% % \begin{proposition}[5-Cycle Spectral Envelope]
% % \label{prop:5-cycle-envelope}
% % The maximal geometric growth rate of any admissible infinite word is bounded by
% % \[
% % \rho_5 =
% % \Bigl(
% % 3^{2}
% % \cdot (123/33)
% % \cdot (341/131)
% % \cdot (329/123)
% % \Bigr)^{1/5}
% % < 3.
% % \]
% % \end{proposition}

% % \begin{lemma}[Cycle Unavoidability]
% % \label{lem:cycle-unavoidability}
% % Every admissible infinite word contains infinitely many occurrences of
% % the 5-cycle (possibly embedded within longer contractive blocks such as
% % 12-cycles). Consequently, no admissible sequence can avoid 5-cycle
% % contributions at bounded intervals.
% % \end{lemma}

% % \begin{remark}
% % Longer contractive cycles (e.g.\ 12-cycles) may delay the appearance of
% % 5-cycles but cannot eliminate them. The grammar forces periodic insertion
% % of 5-cycle structure once sufficient combinatorial slack accumulates.
% % \end{remark}

% % % ------------------------------------------------------------
% % \subsection{Subcritical Growth}

% % \begin{theorem}[Subcritical Symbolic Growth]
% % \label{thm:subcritical-growth-app}
% % Let $\widetilde A(i)$ denote the number of admissible configurations at
% % level $i$. Then
% % \[
% % \limsup_{i \to \infty} \widetilde A(i)^{1/i} \le \rho_5 < 3.
% % \]
% % \end{theorem}

% % \begin{proof}
% % Every admissible infinite word decomposes into concatenations of cycles,
% % each contributing multiplicative growth bounded by Lemma
% % \ref{lem:local-growth-bounds}.

% % By Proposition~\ref{prop:5-cycle-envelope}, the maximal asymptotic growth
% % is achieved by repetition of the 5-cycle, yielding spectral radius
% % $\rho_5$.

% % Since all other cycles are either contractive or sub-maximal, the limsup
% % growth rate is bounded above by $\rho_5$.
% % \end{proof}

% % % ------------------------------------------------------------
% % \subsection{Consequences for Numerator Growth}

% % \begin{corollary}[Dominance by Dyadic Scaling]
% % \label{cor:dyadic-dominance}
% % Since $\rho_5 < 3$ and $2^{A020914(i)} \sim 3^i$, it follows that
% % \[
% % \frac{\widetilde A(i)}{2^{A020914(i)}} \to 0.
% % \]
% % \end{corollary}

% % \begin{remark}[Structural Role]
% % This appendix provides the sole combinatorial input required for
% % Theorem~\ref{thm:maximal-growth}. The main text uses only the existence
% % of a subcritical spectral bound $\rho_5 < 3$ and does not depend on the
% % individual local constants or the internal structure of the cycles.
% % \end{remark}

% % \begin{remark}[Historical Origin]
% % The horizontal-sum decomposition is inspired by the combinatorial
% % summation framework of Winkler~\cite{Winkler2017}, though the present
% % orientation and cycle extraction are adapted specifically to the
% % constraints of the increment grammar governing \OEIS{A022921}.
% % \end{remark}

% % \begin{remark}
% % This appendix contains only structural bounds. All constants are certified
% % in Appendix~\ref{app:local-growth-certification}.
% % \end{remark}

% % % ============================================================
% % \newpage
% % \section{Local Growth Certification}
% % \label{app:local-growth-certification}

% % This appendix provides the explicit finite combinatorial calculations
% % underlying the local growth bounds used in Section~\ref{app:extremal-growth-envelope}.
% % Its role is to certify that all admissible configurations generated by the
% % increment grammar \OEIS{A022921} admit a uniform subcritical growth bound.

% % % ------------------------------------------------------------
% % \subsection{Local Pattern Structure}

% % The dynamics decompose into finitely many local patterns built from single
% % ($s$) and double ($d$) increment steps. Each pattern contributes a
% % multiplicative factor to the total row sum in the horizontal-sum dynamics.

% % The admissible local patterns are:
% % \[
% % (s,d), \quad (d,d), \quad (s,d,s), \quad (d,d,s).
% % \]

% % % ------------------------------------------------------------
% % \subsection{Explicit Growth Bounds}

% % Each pattern admits a finite maximal growth factor, obtained by exhaustive
% % placement of delayed doubling contributions within the horizontal-sum recurrence:

% % \[
% % \begin{array}{c|c}
% % \text{Pattern} & \text{Max Growth Factor} \\
% % \hline
% % s,d & 3 \\
% % d,d & 123/33 \\
% % s,d,s & 341/131 \\
% % d,d,s & 329/123
% % \end{array}
% % \]

% % \begin{remark}
% % These constants arise from finite enumeration of the horizontal-sum recurrence
% % \[
% % a_i(n) = a_i(n-1) + a_{i-1}(n),
% % \]
% % with delayed doubling producing boundary duplication effects.

% % No asymptotic assumptions are used; all values are obtained by exhaustive local evaluation.
% % \end{remark}

% % % ------------------------------------------------------------
% % \subsection{Cycle Decomposition}

% % Admissible infinite words under the grammar decompose into concatenations
% % of finite cycles. Among these, the minimal nontrivial expansive structure is
% % a 5-cycle consisting of the admissible pattern frequencies.

% % The maximal multiplicative growth over one 5-cycle is therefore
% % \[
% % \rho_5 =
% % \Bigl(
% % 3^2 \cdot (123/33) \cdot (341/131) \cdot (329/123)
% % \Bigr)^{1/5}
% % < 3.
% % \]

% % \begin{proposition}
% % All admissible infinite words under \OEIS{A022921} satisfy
% % \[
% % \limsup_{i \to \infty} \widetilde A(i)^{1/i} \le \rho_5 < 3.
% % \]
% % \end{proposition}

% % \begin{proof}
% % Every admissible word decomposes into cycles whose local growth is bounded
% % by the values above. The maximal asymptotic growth is achieved by repetition
% % of the extremal 5-cycle, hence bounded by $\rho_5$.
% % \end{proof}

% % % ------------------------------------------------------------
% % \subsection{Consequence for Global Growth}

% % Since all admissible configurations are generated by concatenation of these
% % cycles, and since each cycle contributes multiplicative growth bounded above
% % by $\rho_5$, the symbolic growth rate of admissible structures is strictly
% % subcritical.

% % \[
% % \frac{\widetilde A(i)}{2^{A020914(i)}} \to 0.
% % \]

% % % ------------------------------------------------------------
% % \subsection{Remark}

% % The detailed row-by-row realizations of extremal growth (including Pascal-type
% % expansions and explicit numerical sequences) are provided only to certify the
% % sharpness of the local bounds. They are not required for the asymptotic
% % argument, which depends solely on the finite pattern structure and cycle
% % decomposition.

% % % ============================================================
% % \newpage
% % \section{Horizontal-Sum Realization of A186009}
% % \label{app:horizontal-sum-table}

% % This appendix provides the explicit horizontal-sum realization of \OEIS{A186009}
% % for reference. It is not used in any asymptotic argument.

% % \begin{table}[htbp]
% % \centering
% % \caption{A186009 Horizontal Sum Formulation}
% % \label{tab:A186009}
% % \vspace{4pt}
% % \begin{tabular}{@{}rrrrrrrrrrr@{}}
% % \toprule
% % $n$ & $a_1(n)$ & $a_2(n)$ & $a_3(n)$ & $a_4(n)$ & $a_5(n)$ &
% % $a_6(n)$ & $a_7(n)$ & $a_8(n)$ & $a_9(n)$ & $\sum a_i(n)$ \\
% % \midrule
% % 1  & 1 &   &   &   &   &   &   &   &   & 1 \\
% % 2  & 1 &   &   &   &   &   &   &   &   & 1 \\
% % 3  & 1 &   &   &   &   &   &   &   &   & 1 \\
% % 4  & 1 & 1 &   &   &   &   &   &   &   & 2 \\
% % 5  & 1 & 2 &   &   &   &   &   &   &   & 3 \\
% % 6  & 1 & 3 & 3 &   &   &   &   &   &   & 7 \\
% % 7  & 1 & 4 & 7 &   &   &   &   &   &   & 12 \\
% % 8  & 1 & 5 & 12 & 12 &   &   &   &   &   & 30 \\
% % 9  & 1 & 6 & 18 & 30 & 30 &   &   &   &   & 85 \\
% % 10 & 1 & 7 & 25 & 55 & 85 &   &   &   &   & 173 \\
% % 11 & 1 & 8 & 33 & 88 & 173 & 173 &   &   &   & 476 \\
% % 12 & 1 & 9 & 42 & 130 & 303 & 476 &   &   &   & 961 \\
% % 13 & 1 & 10 & 52 & 182 & 485 & 961 & 961 &   &   & 2652 \\
% % 14 & 1 & 11 & 63 & 245 & 730 & 1691 & 2652 & 2652 &   & 8045 \\
% % 15 & 1 & 12 & 75 & 320 & 1050 & 2741 & 5393 & 8045 &   & 17637 \\
% % 16 & 1 & 13 & 88 & 408 & 1458 & 4199 & 9592 & 17637 & 17637 & 51033 \\
% % 17 & 1 & 14 & 102 & 510 & 1968 & 6167 & 15759 & 33396 & 51033 & 108950 \\
% % \bottomrule
% % \end{tabular}
% % \end{table}

% % The sequence \OEIS{A186009} arises from the Pascal-type recurrence
% % associated with \OEIS{A100982} as per Winkler~\cite{Winkler2017}:
% % \[
% % a_i(n) = a_i(n-1) + a_{i-1}(n), \quad i,n \in \mathbb{N}.
% % \]

% % \OEIS{A186009} differs only by an initial shift and indexing convention.

% % The single/double increment structure in \OEIS{A022921} corresponds to a
% % boundary modification of this recurrence, inducing the local pattern
% % decomposition used in Appendix~\ref{app:extremal-growth-envelope}.

% % \begin{remark}
% % This appendix provides concrete numerical realization only; all structural
% % and asymptotic results are contained in Appendices ~\ref{app:extremal-growth-envelope}
% % and ~\ref{app:local-growth-certification}.
% % \end{remark}




% % ============================================================
% \section{Appendix: Division Words, Affine Structure, and Extremal Growth}
% \label{app:spine}

% This appendix collects the structural and combinatorial components underlying:
% (i) division-word encoding of Collatz trajectories,
% (ii) their affine representation,
% (iii) admissibility under the A020914 grammar,
% and (iv) the induced extremal growth bounds governing \OEIS{A186009}.

% None of the constructions below are required for the main density argument
% beyond the existence of a subcritical spectral bound.

% % ============================================================
% \subsection{Division Words and Affine Encoding}
% \label{app:spine-affine}

% A division word of length $m$ is a tuple
% \[
% \omega = (d_0,\dots,d_{m-1}), \qquad d_i \in \mathbb{N},
% \]
% recording the 2-adic valuations after each odd step of the fully accelerated map.

% Let
% \[
% D(\omega) = \sum_{i=0}^{m-1} d_i.
% \]

% Each word induces an affine transformation
% \[
% F_\omega(x) = \frac{3^m x + c(\omega)}{2^{D(\omega)}}.
% \]

% \subsubsection*{Canonical constant}

% The constant $c(\omega)$ is uniquely determined modulo $2^{D(\omega)}$ by
% \[
% 3^m x + c(\omega) \equiv 0 \pmod{2^{D(\omega)}}.
% \]

% A canonical representative is chosen by
% \[
% 0 \le c(\omega) < 2^{D(\omega)}.
% \]

% This choice depends only on $\omega$ and not on the representative $x$.

% \subsubsection*{Residue class structure}

% Since $3^m$ is invertible modulo $2^{D(\omega)}$, admissible inputs satisfy
% \[
% x \equiv -c(\omega)\,(3^m)^{-1} \pmod{2^{D(\omega)}}.
% \]

% Thus each division word corresponds to a full arithmetic progression
% (cylinder) of initial conditions.

% % ============================================================
% \subsection{Admissibility and Grammar Constraints}
% \label{app:spine-grammar}

% Let
% \[
% H(i) = \sum_{j=0}^i d_j
% \]
% be the prefix sum.

% Admissible division words of length $m$ satisfy:
% \[
% H(i) < A020914(i+1) \quad (i < m-1),
% \qquad
% H(m-1) = A020914(m).
% \]

% These inequalities define a finite prefix-constrained grammar generating all admissible words.

% \subsubsection*{Enumeration principle}

% Admissible words are obtained by bounded backtracking over $(d_0,\dots,d_{m-1})$
% subject to the prefix constraints above.

% Each admissible word corresponds bijectively to a residue class cylinder
% in the integer lattice.

% % ============================================================
% \subsection{Horizontal-Sum Dynamics and Extremal Growth}
% \label{app:spine-growth}

% The counting sequence \OEIS{A186009} is generated via a horizontal-sum recurrence
% \[
% a_i(n) = a_i(n-1) + a_{i-1}(n),
% \]
% with occasional boundary duplication governed by the A022921 grammar.

% This induces a finite set of local patterns:
% \[
% (s,d), \quad (d,d), \quad (s,d,s), \quad (d,d,s).
% \]

% Each pattern contributes a bounded multiplicative factor to row growth.

% \subsubsection*{Local growth bounds}

% \[
% \begin{array}{c|c}
% \text{Pattern} & \text{Max factor} \\
% \hline
% s,d & 3 \\
% d,d & 123/33 \\
% s,d,s & 341/131 \\
% d,d,s & 329/123
% \end{array}
% \]

% These bounds arise from finite extremal realizations of the recurrence and
% are independent of global row index.

% \subsubsection*{Cycle structure}

% Admissible words decompose into cycles, with the dominant configuration being
% a 5-cycle containing the above pattern frequencies.

% The maximal geometric growth over one cycle is
% \[
% \rho_5 =
% \Bigl(
% 3^2 \cdot (123/33) \cdot (341/131) \cdot (329/123)
% \Bigr)^{1/5}
% < 3.
% \]
% This constant provides the upper spectral envelope used in
% Section~\ref{sec:spectral-bracketing}.

% \subsubsection*{Subcritical growth}

% Let $\widetilde A(i)$ denote admissible counts. Then
% \[
% \limsup_{i\to\infty} \widetilde A(i)^{1/i} \le \rho_5 < 3.
% \]

% Since the dyadic scaling satisfies $2^{A020914(i)} \sim 3^i$, we obtain
% \[
% \frac{\widetilde A(i)}{2^{A020914(i)}} \to 0.
% \]

% Thus admissible cylinders occupy asymptotically vanishing density.

% % ============================================================
% \subsection{Enumeration of Admissible Division Words}
% \label{app:spine-enumeration}

% For fixed $m$, admissible words are those satisfying the prefix constraints
% from Section B.

% Let $A(i+1) = A020914(i+1)$. The enumeration proceeds by constructing
% $(d_0,\dots,d_{m-1})$ such that:
% \[
% H(i) \le A(i+1)-1 \ (i<m-1), \qquad H(m-1)=A(m).
% \]

% \subsubsection*{Algorithmic description}

% The search proceeds lexicographically with backtracking under the constraint:
% \[
% d_i \ge 1, \qquad H(i) < A(i+1).
% \]

% This yields a finite set of admissible words for each $m$.

% \subsubsection*{Example ($m=5$)}

% For $A020914(1{:}5) = (2,4,5,7,8)$, admissible words are exactly those
% satisfying the prefix bounds and summing to $8$ at the terminal index.

% This produces 7 admissible division words, corresponding to
% $A186009(6)=7$.

% % ============================================================
% \subsection{Structural Summary}

% The combined structure is:
% \[
% \text{division word}
% \;\Rightarrow\;
% \text{affine cylinder}
% \;\Rightarrow\;
% \text{prefix grammar constraint}
% \;\Rightarrow\;
% \text{cycle decomposition}
% \;\Rightarrow\;
% \text{subcritical growth}.
% \]

% This pipeline is the only role these constructions play in the main argument:
% they provide a finite combinatorial model whose spectral radius is strictly
% less than the dyadic scaling rate.




% % ============================================================
% \newpage
% \section{A Binomial Lower Comparison for A186009}
% \label{app:A047749-lower-bound}

% \subsection*{Structural Role}

% This sequence provides a \emph{lower-envelope comparison model} for the
% horizontal-sum dynamics underlying \OEIS{A186009}. It corresponds to the
% sub-grammar in which consecutive doubling events are forbidden, producing the
% maximal contractive restriction consistent with the base Pascal-type recursion.
% This sequence provides the lower spectral envelope used in
% Section~\ref{sec:spectral-bracketing}.

% The sequence \OEIS{A047749} admits a horizontal-sum construction analogous to that
% of \OEIS{A186009}, but with a stricter terminal rule: its doubling mechanism is
% governed by \OEIS{A000034}$(n-2)$ rather than \OEIS{A022921}. As a result,
% \OEIS{A047749} forbids \emph{consecutive} doublings, while such events occur
% periodically in \OEIS{A186009}. This structural restriction yields a natural
% comparison sequence.

% \subsection*{Pointwise comparison}

% Because \OEIS{A047749} evolves under the same Pascal-type recursion but with a
% subset of the admissible doubling events, its row sums form a pointwise lower
% bound (up to index shift):
% \[
% A047749(n) \le A186009(n+1), \qquad n \ge 0.
% \]
% The two sequences agree initially, but diverge once the first consecutive
% doubling occurs in \OEIS{A186009}; thereafter the inequality is strict.

% \subsection*{Closed form for \OEIS{A047749}}

% Unlike \OEIS{A186009}, the restricted doubling grammar of \OEIS{A047749}
% permits an exact description. Its terms split naturally into even and odd
% indices:

% \paragraph{Even indices ($n=2m$).}
% \[
% a(n) = \frac{1}{2m+1}\binom{3m}{m}.
% \]

% \paragraph{Odd indices ($n=2m+1$).}
% \[
% a(n) = \frac{1}{2m+1}\binom{3m+1}{m+1}.
% \]

% Equivalently, using Gamma functions,
% \[
% a(2m)=\frac{\Gamma(3m+1)}{\Gamma(2m+2)\Gamma(m+1)}, \qquad
% a(2m+1)=\frac{\Gamma(3m+2)}{\Gamma(2m+2)\Gamma(m+2)}.
% \]

% \subsection*{Asymptotic growth}

% Let $\psi(x)=\Gamma'(x)/\Gamma(x)$ be the digamma function. Differentiating the
% logarithm of the Gamma representations gives

% \[
% \frac{d}{dm}\ln a(2m)
%   = 3\psi(3m+1)-2\psi(2m+2)-\psi(m+1),
% \]
% \[
% \frac{d}{dm}\ln a(2m+1)
%   = 3\psi(3m+2)-2\psi(2m+2)-\psi(m+2).
% \]

% Using $\psi(x)=\ln x + O(1/x)$ as $x\to\infty$,
% both expressions have the same limit:
% \[
% \lim_{m\to\infty} \frac{d}{dm}\ln a(n)
% = \ln\!\left(\frac{3^3}{2^2}\right)
% = \ln\!\left(\frac{27}{4}\right).
% \]

% Thus
% \[
% \lim_{n\to\infty}\frac{a(n+2)}{a(n)}=\frac{27}{4},
% \qquad
% a(n) \sim C\left(\sqrt{\frac{27}{4}}\right)^n.
% \]

% \subsection*{Implication for \OEIS{A186009}}

% Since \OEIS{A047749} is generated by the maximal admissible restriction of the
% A186009 grammar that forbids consecutive doubling, it serves as a strict
% lower-envelope model for symbolic growth.

% Its exponential rate
% \[
% \frac{3\sqrt{3}}{2}
% \]
% therefore provides a universal baseline growth scale for the unrestricted system.
% The additional admissible doubling configurations in \OEIS{A186009} strictly
% increase the available combinatorial branching, placing its growth strictly above
% this baseline.
\newpage
\appendix
\section{Division Words, Grammar Structure, and Spectral Growth}
\label{app:spine}

This appendix collects the structural framework underlying the division-word encoding,
its affine cylinder representation, the admissibility grammar governed by \OEIS{A020914},
and the resulting growth theory for \OEIS{A186009}. The presentation is organized in three layers:
encoding, grammar, and growth.

% ============================================================
\subsection{Affine Encoding of Division Words}
\label{app:spine-affine}

A division word of length $m \ge 1$ is a tuple
\[
\omega = (d_0, d_1, \dots, d_{m-1}), \qquad d_i \in \mathbb{N},
\]
with total exponent
\[
D(\omega) = \sum_{i=0}^{m-1} d_i.
\]

Each word induces an affine transformation
\[
F_\omega(x) = \frac{3^m x + c(\omega)}{2^{D(\omega)}}.
\]

The constant $c(\omega)$ is defined modulo $2^{D(\omega)}$ by
\[
3^m x + c(\omega) \equiv 0 \pmod{2^{D(\omega)}},
\qquad 0 \le c(\omega) < 2^{D(\omega)}.
\]

Let
\[
S_j = \sum_{i=0}^{j-1} d_i, \quad S_0 = 0.
\]

Then
\[
c(\omega)
=
\sum_{j=0}^{m-1}
3^{m-1-j} 2^{S_j}.
\]

Since $3^m$ is invertible modulo $2^{D(\omega)}$, admissible inputs satisfy
\[
x \equiv -c(\omega)(3^m)^{-1} \pmod{2^{D(\omega)}}.
\]

Thus each word defines a cylinder
\[
\mathcal{C}_\omega = \{ r_\omega + k2^{D(\omega)} : k \ge 0 \}.
\]

% ============================================================
\subsection{Grammar Structure and Enumeration}
\label{app:spine-grammar}

Let
\[
H(i) = \sum_{j=0}^i d_j,
\qquad A(i) = \text{\OEIS{A020914}}(i).
\]

A division word is admissible iff
\[
H(i) < A(i+1) \quad (i < m-1),
\qquad
H(m-1) = A(m).
\]

Admissible words are generated by bounded backtracking over
\[
d_i \in \mathbb{N}_{\ge 1}
\]
subject to the prefix constraints.

For $m=5$, with $A(1{:}5) = (2,4,5,7,8)$, there are 7 admissible words,
matching $\text{\OEIS{A186009}}(6)=7$.

Each word corresponds bijectively to a cylinder:
\[
\text{division word} \;\leftrightarrow\; \text{cylinder} \;\leftrightarrow\; \text{affine progression}.
\]

% ============================================================
\subsection{Horizontal-Sum Dynamics and Growth}
\label{app:spine-growth}

The sequence \OEIS{A186009} arises from
\[
a_i(n) = a_i(n-1) + a_{i-1}(n),
\]
with boundary modifications governed by A022921.

This induces local patterns:
\[
(s,d), \quad (d,d), \quad (s,d,s), \quad (d,d,s).
\]

\[
\begin{array}{c|c}
\text{Pattern} & \text{Factor} \\
\hline
(s,d) & 3 \\
(d,d) & 123/33 \\
(s,d,s) & 341/131 \\
(d,d,s) & 329/123
\end{array}
\]

Admissible sequences decompose into cycles. The extremal 5-cycle gives
\[
\rho_5 =
\left(
3^2 \cdot \frac{123}{33} \cdot \frac{341}{131} \cdot \frac{329}{123}
\right)^{1/5}
< 3.
\]

Thus
\[
\limsup_{n\to\infty} A(n)^{1/n} \le \rho_5.
\]

Since $2^{A020914(n)} \sim 3^n$,
\[
\frac{A(n)}{2^{A020914(n)}} \to 0.
\]

% ------------------------------------------------------------
\subsubsection*{Lower spectral envelope (A047749)}
\label{app:A047749-lower-bound}

\OEIS{A047749} corresponds to the restricted grammar forbidding consecutive doubling.

Pointwise:
\[
A047749(n) \le A186009(n+1).
\]

Closed form:
\[
a(2m) = \frac{1}{2m+1}\binom{3m}{m}, \quad
a(2m+1) = \frac{1}{2m+1}\binom{3m+1}{m+1}.
\]

Asymptotic growth:
\[
a(n) \sim C\left(\frac{3\sqrt{3}}{2}\right)^n.
\]

Hence lower rate:
\[
\lambda_{\min} = \frac{3\sqrt{3}}{2}.
\]

% ------------------------------------------------------------
\subsubsection*{Spectral sandwich}

\[
\left(\frac{3\sqrt{3}}{2}\right)^n
\;\lesssim\;
A(n)
\;\ll\;
2^{A020914(n)}.
\]

\[
\text{division words}
\Rightarrow
\text{cylinders}
\Rightarrow
\text{grammar}
\Rightarrow
\text{cycles}
\Rightarrow
\text{spectral bounds}.
\]


% ============================================================
\clearpage
\bibliographystyle{unsrtnat}
\bibliography{src/latex/references}

\end{document}